% v20: complete all-negative branch obstruction for the book low triangle.
% Depends on bk19:canonical, bk19:triangle, bk19:finite-offset, bk19:H-bound.
\section{Book の全負枝低三角形の排除}
\label{bk20:negative-section}

前節の固定 core 有限完成を用いるだけでは，固定 core 自身が動く無限集合を
排除したことにはならない．本節では，全ての二価 cores が負の枝に属する場合に，
低 node の比較から reverse parameters の一様上界を導く．
scale と高分母には探索上限を置かず，非零の因子積から全て復元する．

\begin{theorem}[全負枝の book 低三角形は存在しない]
\label{bk20:no-all-negative}
Book support $G_{5,13}$ の三つの二価 cores が，全て負の枝
$S_i=T_i=w_i<r_i$ に属すると仮定する．
正の整数局所 charts，奇数高分母，三本の低辺の非負 node，
および対向低 $D$ の相殺を同時に満たす低三角形は存在しない．
高中心 cores $0,1$ の完成に関する条件は不要である．
\end{theorem}

\begin{proof}
有向低三角形の順に二価 cores を $i=1,2,3$ と呼び，添字を三を法として読む．
これは高中心 $0,1$ の実際の標識とは別の巡回添字である．
$H_i=r_i-w_i$ と置くと，式~\eqref{bk19:H-bound} により
\begin{equation}\label{bk20:negative-range}
 w_i\ge2,\quad H_i\ge2,\qquad
 \frac{w_i^2}{2}<H_i(H_i+1)<2w_i^2
\end{equation}
である．補題~\ref{bk19:triangle} により canonical 高 ray の向き
$\sigma_i$ は一定ではない．高中心の交換は有向低三角形の向きを反転し，
全ての $\sigma_i$ を交換する．従って，必要なら高中心を交換し巡回添字を選び直すと
\[
 (\sigma_1,\sigma_2,\sigma_3)=(+,+,-)
\]
としてよい．

全て $L_i=1$ なので，命題~\ref{bk19:offsets} の有限 offset は
\begin{equation}\label{bk20:negative-offsets}
 w_2=r_1+t_2,\qquad r_3=r_2+t_3,\qquad w_1=w_3+t_1,
\end{equation}
\[
 -1\le t_2\le2,\qquad -2\le t_3\le2,\qquad -1\le t_1\le1
\]
となる．ここで node 条件は三つの高分母そのものを使わず，reverse entries のみを使う．

\paragraph{第一段階：証明による一様上界．}
式~\eqref{bk20:negative-range} より
\[
 H+\frac12>\frac{w}{\sqrt2}>\frac7{10}w,
 \qquad H<\sqrt2w<\frac{10}{7}w
\]
なので
\begin{equation}\label{bk20:r-bounds}
 \frac{17}{10}w-\frac12<r<\frac{17}{7}w
\end{equation}
を得る．式~\eqref{bk20:negative-offsets} と合わせると
\[
 w_2>\frac{17}{10}w_1-\frac32,
\]
\[
 r_3>\frac{17}{10}w_2-\frac52
      >\frac{289}{100}w_1-\frac{101}{20},
 \qquad
 r_3<\frac{17}{7}w_3\le\frac{17}{7}(w_1+1).
\]
従って
\begin{equation}\label{bk20:proved-bound}
 \frac{323}{700}w_1<\frac{1047}{140},\qquad
 w_1<\frac{5235}{323}<17.
\end{equation}
整数性から $w_1\le16$，$w_3\le17$ であり，
$w_2=r_1+t_2<\frac{17}{7}\cdot16+2<41$ から $w_2\le40$ である．
これらは全て仮定から導いた上界である．

\paragraph{第二段階：固定 reverse data の全 scale．}
負枝の canonical 表示で
\[
 c=\frac{sH}{r},\qquad e=\frac{z(H+1)}{r+1},\qquad
 E=1-\frac{w^2}{H(H+1)}=\frac MN
\]
と置き，$M,N$ を互いに素，$N>0$ とする．
式~\eqref{ml12:E-equation} と式~\eqref{kb15:E-bounds} は
\[
 Ece=e-c+1,\qquad c\ge2,\ e\ge1,\quad -1<E<\frac12,\quad E\ne0
\]
を与える．従って
\begin{equation}\label{bk20:complete-factorization}
 (N-Mc)(N+Me)=N(N-M)>0.
\end{equation}
両因子は正である．実際，$Ec=1-(c-1)/e<1$ と
$c(1+Ee)=e+1>0$ を用いればよい．
よって $d$ を $N(N-M)$ の全ての正の約数として
\[
 c=\frac{N-d}{M},\qquad
 e=\frac{N(N-M)/d-N}{M},\qquad
 s=\frac{rc}{H},\qquad z=\frac{(r+1)e}{H+1}
\]
を復元し，正整数性と奇数分母 $sw+1,zw+1$ を確認すれば，全ての scale が得られる．
式~\eqref{bk20:complete-factorization} の右辺は非零なので，共鳴直線の見落としはない．
$s,z$ または高分母に上限を加えていない．

\paragraph{第三段階：残る有限データの厳密検査．}
$2\le w_1\le16$，式~\eqref{bk20:negative-range}，および
式~\eqref{bk20:negative-offsets} を満たす reverse data は九十一型である．
第二段階の全 scale 復元が三つの cores 全てで空でないものは二十五型であり，
scale を展開した三模型の組は六十四個となる．
その全てで，最初の低辺の $D$ 相殺が既に破れる．

次の表は，第一 core の canonical $(r,s,z,S,T)$ ごとに，
可能な第二 core との残差 $D_C^{(1)}+D_D^{(2)}$ をまとめたものである．
異なる第一・第二 core の組は二十六組であり，表のいずれの集合にも零はない．
\begin{center}
\small
\begin{tabular}{ll}
第一 core の canonical parameters & 第一低辺の $D$ 残差の全値\\
\hline
$(4,4,5,2,2)$ & $-20,-9,-4,4,7,16$\\
$(5,10,2,3,3)$ & $-40,-25,-9,-1,11,24$\\
$(6,6,14,3,3)$ & $-32,32$\\
$(7,7,2,4,4)$ & $-57,-20,-9,3,12,43$\\
$(7,21,4,4,4)$ & $-64,-27,-16,-4,5,36$\\
\end{tabular}
\end{center}

この有限検査は \nolinkurl{verify_book_all_negative_v20.py} に実装し，
九十一型の reverse data と全 scale，六十四個の node と $D$ 証明書を
\nolinkurl{book_all_negative_verification_v20.json} に保存した．
行和，pair products，明示的な正 weights，三つの node の非負性と総数五は
元の行列から再計算している．
数学的に導いた有限集合上の整数演算であり，有限範囲の探索結果から
無上界の結論へ外挿する段階は含まれない．
以上で全負枝を排除した．
\end{proof}

\begin{corollary}[Book には正負両枝が必要]
\label{bk20:mixed-necessary}
仮定を満たす book 低三角形が存在するなら，三つの二価 cores の中に
正の枝と負の枝がともに現れる．従って負の枝は一個または二個である．
\end{corollary}

\begin{proof}
全正枝は補題~\ref{bk19:triangle}，全負枝は
定理~\ref{bk20:no-all-negative} により排除される．
\end{proof}

この系は book support 全体の排除ではない．
正負混合の低三角形，高中心 cores の完成，および残る支持の一般分母剛性は
引き続き別の問題である．
